\documentclass[9pt,a4paper]{article}

\usepackage[a4paper,top=12mm,bottom=12mm,left=15mm,right=15mm]{geometry}
\usepackage{fontspec}
\usepackage{xcolor}
\usepackage{tabularx}
\usepackage{enumitem}
\usepackage[most]{tcolorbox}
\usepackage{microtype}
\usepackage{array}

\IfFontExistsTF{Arial}{
  \setmainfont{Arial}
  \setsansfont{Arial}
}{
  \setmainfont{TeX Gyre Heros}
  \setsansfont{TeX Gyre Heros}
}

\definecolor{AIEPRed}{HTML}{D62027}
\definecolor{AIEPNavy}{HTML}{082B5C}
\definecolor{AIEPInk}{HTML}{182B3A}
\definecolor{AIEPMuted}{HTML}{5F6B7A}
\definecolor{AIEPLine}{HTML}{D8DEE6}
\definecolor{AIEPSoft}{HTML}{F5F2EC}

\pagestyle{empty}
\setlength{\parindent}{0pt}
\setlength{\parskip}{3pt}
\setlist[itemize]{leftmargin=*,topsep=1pt,itemsep=1pt,parsep=0pt}
\setlist[enumerate]{leftmargin=*,topsep=1pt,itemsep=1pt,parsep=0pt}
\renewcommand{\arraystretch}{1.08}

\newtcolorbox{softbox}[1]{
  colback=AIEPSoft,
  colframe=AIEPLine,
  boxrule=0.5pt,
  arc=1mm,
  left=2.5mm,right=2.5mm,top=1.7mm,bottom=1.7mm,
  before skip=3pt,after skip=3pt,
  title={\scriptsize\bfseries\MakeUppercase{#1}},
  coltitle=AIEPNavy,
  attach title to upper={\par\vspace{0.5mm}}
}

\newtcolorbox{accentbox}[1]{
  colback=white,
  colframe=AIEPRed,
  boxrule=0.65pt,
  arc=1mm,
  left=2.5mm,right=2.5mm,top=1.7mm,bottom=1.7mm,
  before skip=3pt,after skip=3pt,
  title={\scriptsize\bfseries\MakeUppercase{#1}},
  coltitle=AIEPRed,
  attach title to upper={\par\vspace{0.5mm}}
}

\begin{document}

{\small\bfseries\textcolor{AIEPRed}{GEOGREEN ESCOLAR · AIEP OSORNO}}\par
\vspace{2mm}
{\Huge\bfseries\color{AIEPNavy}Léeme — Mentoría 2\par}
\vspace{1mm}
{\Large\bfseries\color{AIEPInk}Solución propuesta y recursos definidos\par}
{\normalsize\color{AIEPMuted}Material base para el equipo docente de Ingeniería, Energía y Tecnología (I/E/T)\par}

\vspace{4mm}
Esta carpeta reúne los materiales para ejecutar la Mentoría 2 de GeoGreen Escolar. Puede utilizarlos directamente o ajustar su mediación al grupo. La organización, los productos de entrada y el entregable final corresponden a la lógica común del programa y deben preservarse.

\begin{softbox}{Qué hace esta mentoría}
La sesión convierte un problema validado y una idea tecnológica inicial en un producto mínimo viable definido. Para ello, cada equipo:
\begin{enumerate}
  \item describe la solución como observar, decidir, responder y mostrar;
  \item define una variable, el sensor o fuente de datos y un umbral fundamentado;
  \item identifica recursos físicos, capa de software y condiciones de seguridad;
  \item acota el alcance y distribuye tareas entre sus seis responsabilidades;
  \item compromete una evidencia concreta para la Mentoría 3.
\end{enumerate}
\end{softbox}

\begin{accentbox}{Condiciones de ejecución}
\begin{tabular}{@{}>{\bfseries\color{AIEPNavy}\raggedright\arraybackslash}p{39mm}@{\hspace{5mm}}>{\raggedright\arraybackslash}p{\dimexpr\linewidth-44mm\relax}@{}}
\strut Duración & \strut 60 minutos por bloque. \\[1mm]
\strut Organización & \strut 5 equipos de 6 estudiantes por bloque. \\[1mm]
\strut Entradas obligatorias & \strut Problema y contexto M1, idea tecnológica T3 y próximo paso comprometido. \\[1mm]
\strut Salida obligatoria & \strut Solución y recursos definidos: regla, software, seguridad, tareas y evidencia para M3. \\[1mm]
\strut Rol de la mentoría & \strut Orientar, contrastar y resguardar la seguridad; la construcción continúa con trabajo del equipo entre sesiones. \\
\end{tabular}
\end{accentbox}

\begin{softbox}{Preparación breve}
Revise la presentación e imprima una ficha por equipo —cinco por bloque—. Solicite los productos de entrada, disponga el banco acotado de sensores y el manual del kit, compruebe el proyector y establezca un punto de revisión física antes de conectar cualquier montaje por USB.
\end{softbox}

\begin{accentbox}{Atención especial}
El problema orienta la tecnología, no al revés. Se prioriza una variable principal y una prueba mínima. Voltaje, polaridad, tierra común, resistencias y montaje se revisan físicamente; ante una duda, el circuito permanece sin energía. Cualquier intervención sobre 220 V requiere supervisión docente.
\end{accentbox}

{\bfseries\color{AIEPNavy}Qué hay en esta carpeta}\par
\begin{itemize}
  \item \textbf{LEEME.pdf} — esta guía de orientación.
  \item \textbf{GeoGreen Escolar - resumen del proyecto.png} y \textbf{GeoGreen Escolar - lógica del proyecto.png} — panorama general y continuidad de talleres, mentorías y productos.
  \item \textbf{Documentos/} — planificación docente completa y ficha imprimible para el trabajo de cada equipo.
  \item \textbf{Presentacion/} — apoyo visual para conducir la sesión en PowerPoint.
\end{itemize}

\begin{tcolorbox}[colback=AIEPNavy,colframe=AIEPNavy,coltext=white,boxrule=0pt,arc=1mm,left=3mm,right=3mm,top=1.8mm,bottom=1.8mm,before skip=3pt]
\textbf{Resultado esperado:} cada equipo sale con un producto mínimo viable acotado, una regla explicable, recursos y seguridad definidos, seis tareas verificables y una evidencia comprometida para el 21 de septiembre.
\end{tcolorbox}

\vfill
{\footnotesize\color{AIEPMuted}GeoGreen Escolar · AIEP Osorno — Ingeniería, Energía y Tecnología (I/E/T) · Desarrollo Social.}

\end{document}
